
\begin{table}[t]
\caption{RQ3: share of each era's candidates by author role (\%), and direction--outcome alignment.}
\label{tab:rq3roles}
\begin{tabular}{lrr}
\toprule
 & BDFL era & SC era \\
\midrule
Community member & 40.5 & 35.9 \\
Core developer & 27.9 & 30.9 \\
Proposal author & 16.4 & 14.4 \\
PEP editor & 8.1 & 5.9 \\
BDFL & 5.1 & 0.0 \\
BDFL / PEP delegate & 1.9 & 0.8 \\
Steering Council member & --- & 12.1 \\
\midrule
Participants & 2,310 & 1,546 \\
Proposals & 494 & 579 \\
Supporting aligned with positive outcome & 73.5 & 62.7 \\
Blocking aligned with negative outcome & 27.5 & 23.8 \\
Revising aligned with undecided outcome & 6.3 & 17.1 \\
\bottomrule
\end{tabular}
\end{table}

\paragraph{Medium or era?} Two thirds of the Steering Council era's candidates come from Discourse and one third from the mailing lists. The two venues have nearly identical mechanism profiles within the era (operational 28.3\% vs.\ 30.0\%, functional 29.0\% vs.\ 28.9\%, authority 9.8\% vs.\ 9.5\%, strategic 5.2\% vs.\ 6.5\%; Discourse carries more ecosystem signals, 19.2\% vs.\ 13.4\%, and fewer compatibility ones, 11.9\% vs.\ 15.1\%), and Steering Council members account for a similar share on both (11.1\% vs.\ 13.3\%). The differences between eras are therefore not an artefact of the change of medium, although the lower unknown-scope share and the username-based identities of Discourse should be kept in mind.
